\begin{trapbox}
\textbf{误区一：左右焦点可以随意配准线。}
\begin{center}\begin{tikzpicture}[font=\small,>=stealth]
\draw[black!30,rounded corners=2pt] (0,0) rectangle (14.4,8.2);
\node[anchor=west,font=\small\bfseries] at (.5,7.72) {配错准线的反例：同一个点、同一个焦点、同一尺度};
\begin{scope}[shift={(3.6,4)},scale=.34]
\draw[blue,thick] (0,0) ellipse (5 and 3);\draw[orange!85!black,very thick] (4,0)--(3,2.4);
\draw[violet,thick] (6.25,-4)--(6.25,4);\draw[violet,very thick] (3,2.4)--(6.25,2.4);
\fill (4,0) circle (2pt) node[left=5pt] {$F_R$};\fill (3,2.4) circle (2pt) node[above=8pt] {$P$};
\fill (6.25,2.4) circle (2pt) node[right] {$H$};
\end{scope}
\node at (3.6,6.6) {右焦点 + 右准线};
\node at (3.6,1.4) {$(13/5)/(13/4)=4/5$};
\begin{scope}[shift={(10.8,4)},scale=.34]
\draw[blue,thick] (0,0) ellipse (5 and 3);\draw[orange!85!black,very thick] (4,0)--(3,2.4);
\draw[violet,thick] (-6.25,-4)--(-6.25,4);\draw[violet,very thick] (3,2.4)--(-6.25,2.4);
\fill (4,0) circle (2pt) node[left=5pt] {$F_R$};\fill (3,2.4) circle (2pt) node[above=8pt] {$P$};
\fill (-6.25,2.4) circle (2pt) node[left] {$H$};
\end{scope}
\node at (10.8,6.6) {右焦点 + 左准线};
\node at (10.8,1.4) {$(13/5)/(37/4)=52/185$};

\node[anchor=west] at (.7,.55) {右图比值不等于 $e=4/5$；应保持 $F_R\leftrightarrow\ell_R$、$F_L\leftrightarrow\ell_L$};
\end{tikzpicture}\end{center}

右焦点应配右准线，左焦点应配左准线。个别点即便偶然给出相同比值，也不能推出整条曲线恒成立。

\textbf{误区二：把 $PF=e\,d$ 反解成 $d=e\,PF$。}
正确为 $d=PF/e$；$0<e<1$ 时准距大于焦半径，可先用大小关系自查。
\end{trapbox}
\clearpage
\begin{trapbox}
\textbf{误区三：焦点在 $y$ 轴仍写竖直准线。}
\begin{center}\begin{tikzpicture}[font=\small,>=stealth]
\draw[black!30,rounded corners=2pt] (0,0) rectangle (14.4,8.2);
\node[anchor=west,font=\small\bfseries] at (.5,7.72) {纵向椭圆 $y^2/25+x^2/9=1$：上焦点配上准线};
\begin{scope}[shift={(6.7,3.6)},scale=.47]
\draw[blue,very thick] (0,0) ellipse (3 and 5);
\draw[violet,thick] (-5,6.25)--(5,6.25) node[right] {$\ell_U:y=25/4$};
\draw[black!35] (0,-5.3)--(0,6.5);
\draw[orange!85!black,very thick] (0,4)--(2.4,3);
\draw[violet,very thick] (2.4,3)--(2.4,6.25);
\fill (0,4) circle (2pt);
\node[anchor=east] at (-4.2,4) {$F_U(0,4)$};
\draw[black!45,->] (-3.8,4)--(0,4);
\fill (2.4,3) circle (2pt) node[right=7pt] {$P(12/5,3)$};
\fill (2.4,6.25) circle (2pt) node[above] {$H$};
\end{scope}

\node[anchor=west] at (.7,.55) {焦点轴为 $y$ 轴，准线为水平线；比值仍为 $(13/5)/(13/4)=4/5$};
\end{tikzpicture}\end{center}

准线始终垂直于焦点轴；纵向椭圆取 $y=\pm a^2/c$。

\textbf{误区四：把 $e=0$ 或 $e=1$ 直接当作本条椭圆。}
本条要求 $0<e<1$。$c=0$ 的圆不能代入 $a^2/c$；对固定的有限焦点、准线，$e=0$ 的距离等式只有点 $F$，
$e=1$ 对应抛物线，均不能套用本条的非退化椭圆方程。
\end{trapbox}
